\section{Main Results}
\label{sec:main-results}

\subsection{Headline Theorems}
\label{sec:headline}

\TW{} delivers four standard headline guarantees in the lifted
public-permutation model of \Cref{sec:games-pp}: all-in-one AE
security \cite{RS06}, its multi-user indistinguishability lift
\cite{BT16}, and the two facets of commitment security---$s$-way
CMTDs-4 committing security \cite{BH22} and context discoverability
\cite{MLGR23}. Committing security is the collision facet and context
discoverability the preimage facet; both follow from the same root tag
and are stated co-equally. We also state tag-projected strengthenings
showing that the root tag alone is the compact commitment target. For
the AE-style notions the ideal side
exposes the BDPV simulator; the committing and discoverability bounds
remain real public-permutation success probabilities.
\Cref{thm:ae,thm:mu-ind,thm:cmtds4,thm:tag-cmtds4,thm:cdy,thm:tag-cdy}
state these bounds; the loss terms $\Lift_c$, $\KeyGuess_k$,
$\RootColl_\tau$, $\LeafColl_\tau$, and
$\RootPre_\tau$
are unpacked in \Cref{sec:loss-terms}, and the random-oracle envelopes
$\Epspriv$, $\Epsintctxt$ in \Cref{sec:thms-pp}.

\begin{theorem}[All-in-One AE Security]
\label{thm:ae}
For every nonce-respecting AE adversary $\mathcal{A}$ against \TW{}
with the resources of \Cref{sec:resources},
\[
  \Advae_\TW(\mathcal{A})
  \;\leq\;
  \Lift_c(\Ntot)
  + \Epspriv(\Nprim)
  + \Epsintctxt(\Nprim, \Nleaf, \Qv).
\]
\end{theorem}

\begin{proof}
Let
\[
  p_0 =
    \Pr[\mathcal{A}^{\Enc_K[\pi],\, \Ver_K[\pi],\, \pi,\, \pi^{-1}}
        \Rightarrow 1],
  \qquad
  p_1 =
    \Pr[\mathcal{A}^{\Rand,\, \Replay,\,
                    \Sim^{\ROflat},\, \Sim^{-1,\ROflat}}
        \Rightarrow 1].
\]
Then $\Advae_\TW(\mathcal{A}) = |p_0 - p_1|$ by
\Cref{sec:games-pp}. By the flat-sponge lift bound
(\Cref{lem:lift-bound}), replacing the real primitive interfaces by the
simulator and replacing the construction by $(\Enc_K^\RO,\Ver_K^\RO)$
costs at most $\Lift_c(\Ntot)$:
\[
  |p_0 - p_0'| \leq \Lift_c(\Ntot),
  \quad
  p_0' =
    \Pr[\mathcal{A}^{\Enc_K^\RO,\,\Ver_K^\RO,\,
                    \Sim^{\ROflat},\,\Sim^{-1,\ROflat}}
        \Rightarrow 1].
\]
Let $\mathcal{B} =
\mathcal{B}_{\mathrm{derived}}(\mathcal{A})$ be the adversary of
\Cref{lem:derived-adversary}. The derived-adversary construction gives
\[
  p_0' =
    \Pr[\mathcal{B}^{\Enc_K^\RO,\,\Ver_K^\RO,\,\ROflat} \Rightarrow 1],
  \qquad
  p_1 =
    \Pr[\mathcal{B}^{\Rand,\,\Replay,\,\ROflat} \Rightarrow 1],
\]
with $\qRO(\mathcal{B}) \leq \Nprim$,
$\nleaf(\mathcal{B}) \leq \Nleaf$, and
$\qv(\mathcal{B}) \leq \Qv$. Therefore
\[
  \Advae_\TW(\mathcal{A})
  \leq \Lift_c(\Ntot) + \Adv^{\mathsf{ae}}_{\RO}(\mathcal{B}).
\]
The all-in-one composition lemma (\Cref{lem:bn-composition}), applied
in the random-oracle model, gives
\[
  \Adv^{\mathsf{ae}}_{\RO}(\mathcal{B})
  \;\leq\;
  \Adv^{\mathsf{priv}}_{\RO}(\mathcal{B}_{\text{priv}})
  + \Adv^{\mathsf{int\text{-}ctxt}}_{\RO}(\mathcal{B}_{\text{int-ctxt}}),
\]
for derived adversaries $\mathcal{B}_{\text{priv}}$ (which runs
$\mathcal{B}$, records the outputs of its encryption-oracle queries,
and answers each verification query by replay-table lookup) and
$\mathcal{B}_{\text{int-ctxt}}$ (which runs $\mathcal{B}$, forwards
encryption and verification queries to its own oracles, and halts on
the first accepting non-replay verification response), both inheriting
$\mathcal{B}$'s $(\qRO, \nleaf, \qv)$. Because $\mathcal{B}_{\text{priv}}$
forwards $\mathcal{B}$'s encryption queries unchanged and $\mathcal{B}$
inherits the nonce-respecting discipline of $\mathcal{A}$, the adversary
$\mathcal{B}_{\text{priv}}$ is itself nonce-respecting, so the privacy
envelope of \Cref{sec:thms-pp} applies.
Substituting the random-oracle envelopes of \Cref{sec:thms-pp} gives
the stated bound.
\end{proof}

\begin{lemma}[Mu-Ind Adjacent-Hybrid Bound]
\label{lem:mu-single-user-envelope}
Let $\mathcal{B}$ be a flat-RO mu-ind adversary in the pre-sampled-key
formulation of \Cref{sec:mu-setting}, with keys
$K_1,\dots,K_D$ sampled at the start of the experiment. Fix
$d \in \{1,\dots,D\}$. Let $H_{d-1}^\star$ and $H_d^\star$ be the
adjacent stopped hybrids that output a fixed bit if any two sampled keys
collide, and otherwise run $\mathcal{B}$ with one shared lazy
$\ROflat$ table as follows: users $e < d$ are answered by
$(\Rand,\Replay)$, users $e > d$ are answered by
$(\Enc_{K_e}^\RO,\Ver_{K_e}^\RO)$, and user $d$ is answered by
$(\Enc_{K_d}^\RO,\Ver_{K_d}^\RO)$ in $H_{d-1}^\star$ and by
$(\Rand,\Replay)$ in $H_d^\star$. If $\nleaf^{(d)}$ and $\qv^{(d)}$
are the user-$d$ resource caps, then
\[
  \bigl|\Pr[\mathcal{B}^{H_{d-1}^\star} \Rightarrow 1]
       - \Pr[\mathcal{B}^{H_d^\star} \Rightarrow 1]\bigr|
  \leq
  2\KeyGuess_k(\qRO(\mathcal{B}))
  + \frac{\nleaf^{(d)}(\nleaf^{(d)}-1)}{2^{\tauleaf+1}}
  + \frac{\qv^{(d)}}{2^{\tauroot}}.
\]
\end{lemma}

\begin{proof}
The fixed-output branch on a key collision is identical in
$H_{d-1}^\star$ and $H_d^\star$, so only the branch with pairwise
distinct keys contributes. On that branch, the two hybrids differ only
in user $d$: users $e < d$ are ideal in both hybrids, and users
$e > d$ are real in both hybrids.

Let $\mathcal{E}_d$ denote the common surrounding environment. It runs
$\mathcal{B}$, answers users $e < d$ by their $(\Rand,\Replay)$
interfaces, answers users $e > d$ by
$\Enc_{K_e}^\RO,\Ver_{K_e}^\RO$, forwards $\mathcal{B}$'s direct
$\ROflat$ queries unchanged, and shares the same lazy $\ROflat$ table
with user $d$. The Bellare--Namprempre insertion is pathwise in this
environment. Insert the interface in which user $d$'s encryption oracle
is still real but user $d$'s verification oracle is replaced by its
replay table. The first hop is the user-$d$ INT-CTXT event; the second
hop is the user-$d$ privacy replacement with the same environment and
replay table. The replay table is determined by user $d$'s encryption
answers and adds no direct $\ROflat$ queries or construction-generated
transcripts.

The privacy and INT-CTXT accounting used in these two hops is local to
user $d$'s keyed transcript language. In the privacy hop, the
one-query replacement of \Cref{lem:priv-one-query-replacement} only
requires that no direct query has read a user-$d$ counted transcript
before $\badkey^{(d)}$, the $\badkey$ event of
\Cref{sec:privacy,sec:authenticity} instantiated at user $d$'s key
$K_d$; any extra lazy-table entries
written by the real user-$d$ encryption computation are unobservable
before that abort. The accepted user-$d$ encryption queries are
nonce-respecting because the mu-ind interface enforces nonce uniqueness
per user. In the INT-CTXT hop, the analysis of \Cref{thm:intctxt-ro}---its
visible-prefix key uniformity (\Cref{lem:int-key-uniformity}), its
leaf-freshness and encryption-first-replay steps, and the
verification-first root-guessing bound
(\Cref{lem:int-verification-first-root-guess})---refers only to user-$d$
hidden-key transcript points, user-$d$ final leaf transcripts, and
user-$d$ non-replay verification submissions.

It remains to justify that the same-RO environment does not enlarge the
direct-query term. The real users $e > d$ may perform
construction-internal $\ROflat$ lookups while simulating
$\Enc_{K_e}^\RO$ and $\Ver_{K_e}^\RO$. Whenever such a lookup lies in a
counted keyed transcript class, \Cref{lem:bridge-decode} decodes its
candidate key as $K_e$; on the non-collision branch this is different
from $K_d$. If the lookup is outside the counted keyed language, it is
irrelevant to $\badkey^{(d)}$. Thus foreign-user
construction lookups neither trigger user $d$'s key-test event nor
pre-read a user-$d$ hidden transcript point. They also do not create
user-$d$ leaf tag collisions or user-$d$ verification attempts, which
are counted by $\nleaf^{(d)}$ and $\qv^{(d)}$.

The key-test contribution is taken in the original pre-sampled-key
experiment, not in a renormalized experiment conditioned on the
non-collision branch. Fix the foreign keys and run the visible-prefix
coupling used in the user-$d$ privacy and INT-CTXT proofs, with
foreign-key transcript addresses left unchanged. Before the first
user-$d$ key hit, the visible prefix is invariant under renaming $K_d$
among keys not already tested and not equal to any fixed foreign key.
For each fresh candidate key tested by one of $\mathcal{B}$'s original
direct $\ROflat$ queries, the event contributing on the stopped branch
is contained in $\{K_d \text{ equals that candidate}\}$; if the
candidate is a foreign key, it contributes nothing on the
non-collision branch. Hence each fresh original direct test costs at
most $2^{-k}$ in the unconditioned key experiment, and at most
$\qRO(\mathcal{B})$ such tests occur.

The proofs of \Cref{thm:priv-ro,thm:intctxt-ro} therefore apply to the
two user-$d$ hops with the same direct-query cap
$\qRO(\mathcal{B})$ and the same user-$d$ construction resources
$\nleaf^{(d)},\qv^{(d)}$. The surrounding environment affects oracle
scheduling and foreign construction-internal table entries, but not the
counted user-$d$ events. The privacy hop contributes
$\KeyGuess_k(\qRO(\mathcal{B}))$, and the INT-CTXT hop contributes
\[
  \KeyGuess_k(\qRO(\mathcal{B}))
  + \frac{\nleaf^{(d)}(\nleaf^{(d)}-1)}{2^{\tauleaf+1}}
  + \frac{\qv^{(d)}}{2^{\tauroot}}.
\]
Combining the two hops gives the stated bound.
\end{proof}

\begin{theorem}[Multi-User Indistinguishability]
\label{thm:mu-ind}
For every mu-ind adversary $\mathcal{A}$ against \TW{} with
execution-uniform user cap $D$ and the global resources of
\Cref{sec:mu-setting},
\[
  \Advmuind_\TW(\mathcal{A})
  \;\leq\;
  \Lift_c(\Ntot)
  + \frac{D(D - 1)}{2^{k + 1}}
  + 2D \cdot \KeyGuess_k(\Nprim)
  + \frac{\Nleaf(\Nleaf - 1)}{2^{\tauleaf + 1}}
  + \frac{\Qv}{2^{\tauroot}}.
\]
\end{theorem}

\begin{proof}
Let $\mathsf{Real}^{\mathsf{mu}}_\TW[\pi]$ denote the $b = 1$
multi-user construction interface of \Cref{sec:mu-setting}, with user
$d$ using $(\Enc_{K_d}[\pi],\Ver_{K_d}[\pi])$, and let
$\mathsf{Ideal}^{\mathsf{mu}}$ denote the $b = 0$ interface with
per-user $(\Rand,\Replay)$. Use the pre-sampled-key formulation of
\Cref{sec:mu-setting}: sample keys $K_1,\dots,K_D$ at the start of the
experiment, and let the $d$-th $\mathsf{New}$ call activate $K_d$.
Unused keys are ignored. Define
\[
  p_0 =
    \Pr[\mathcal{A}^{\mathsf{Real}^{\mathsf{mu}}_\TW[\pi],\,
                    \pi,\,\pi^{-1}} \Rightarrow 1],
  \qquad
  p_1 =
    \Pr[\mathcal{A}^{\mathsf{Ideal}^{\mathsf{mu}},\,
                    \Sim^{\ROflat},\,\Sim^{-1,\ROflat}} \Rightarrow 1].
\]
Then $\Advmuind_\TW(\mathcal{A}) = |p_0 - p_1|$ by
\Cref{sec:mu-ind,sec:mu-setting}. By the flat-sponge lift bound
(\Cref{lem:lift-bound}), replacing $(\pi,\pi^{-1})$ on the real side by
$(\Sim^{\ROflat},\Sim^{-1,\ROflat})$ and each real user construction
oracle by its random-oracle form costs at most $\Lift_c(\Ntot)$:
\[
  |p_0 - p_0'| \leq \Lift_c(\Ntot),
  \quad
  p_0' =
    \Pr[\mathcal{A}^{\mathsf{Real}^{\mathsf{mu}}_{\RO},\,
                    \Sim^{\ROflat},\,\Sim^{-1,\ROflat}} \Rightarrow 1],
\]
where $\mathsf{Real}^{\mathsf{mu}}_{\RO}$ uses
$(\Enc_{K_d}^{\RO},\Ver_{K_d}^{\RO})$ for every created user $d$.
Let $\mathcal{B} =
\mathcal{B}_{\mathrm{derived}}(\mathcal{A})$ be the adversary of
\Cref{lem:derived-adversary}. The derived-adversary construction gives
\[
  p_0' =
    \Pr[\mathcal{B}^{\mathsf{Real}^{\mathsf{mu}}_{\RO},\,\ROflat}
        \Rightarrow 1],
  \qquad
  p_1 =
    \Pr[\mathcal{B}^{\mathsf{Ideal}^{\mathsf{mu}},\,\ROflat}
        \Rightarrow 1],
\]
with $\qRO(\mathcal{B}) \leq \Nprim$ and the same per-user
construction-side transcript caps. Therefore
\[
  \Advmuind_\TW(\mathcal{A})
  \leq \Lift_c(\Ntot)
    + \Adv^{\mathsf{mu\text{-}ind}}_{\RO}(\mathcal{B}).
\]
Let $\mathsf{coll}$ denote the event that two of the pre-sampled keys
$K_1,\dots,K_D$ collide;
$\Pr[\mathsf{coll}] \leq \binom{D}{2}/2^k = D(D-1)/2^{k+1}$ by union
bound. In the random-oracle model, use a standard hybrid over the $D$
users: $H_d$ answers users $1,\dots,d$ with $(\Rand,\Replay)$ and
users $d+1,\dots,D$ with the real construction interface, so $H_0$ is
the real mu-ind game and $H_D$ is the ideal one. Let $H_d^\star$ be
$H_d$ with a fixed output bit on $\mathsf{coll}$. Since
$\mathsf{coll}$ has the same probability in all hybrids,
\[
  \Adv^{\mathsf{mu\text{-}ind}}_{\RO}(\mathcal{B})
  \leq
  \Pr[\mathsf{coll}]
  +
  \bigl|\Pr[\mathcal{B}^{H_0^\star} \Rightarrow 1]
       - \Pr[\mathcal{B}^{H_D^\star} \Rightarrow 1]\bigr|.
\]
By the triangle inequality, it remains to bound each adjacent stopped
hybrid.
For each $d$, \Cref{lem:mu-single-user-envelope} gives the adjacent
hybrid bound
\[
  \bigl|\Pr[\mathcal{B}^{H_{d-1}^\star} \Rightarrow 1]
       - \Pr[\mathcal{B}^{H_d^\star} \Rightarrow 1]\bigr|
  \leq
  2\KeyGuess_k(\qRO(\mathcal{B}))
  + \frac{\nleaf^{(d)}(\nleaf^{(d)}-1)}{2^{\tauleaf+1}}
  + \frac{\qv^{(d)}}{2^{\tauroot}}.
\]

Using $\qRO(\mathcal{B}) \leq \Nprim$, each hybrid step costs
$2\KeyGuess_k(\Nprim)$ plus the per-user leaf-collision and
verification-tag terms. Summing across $d = 1,\dots,D$ gives
$2D \KeyGuess_k(\Nprim)$ for the key-guess contribution. The
leaf-collision contribution satisfies
\[
  \sum_d \frac{\nleaf^{(d)}(\nleaf^{(d)} - 1)}{2^{\tauleaf + 1}}
  = \frac{1}{2^{\tauleaf}} \sum_d \binom{\nleaf^{(d)}}{2}
  \leq \frac{1}{2^{\tauleaf}} \binom{\sum_d \nleaf^{(d)}}{2}
  \leq \frac{\Nleaf(\Nleaf - 1)}{2^{\tauleaf + 1}},
\]
the first inequality because the unordered pairs drawn within the
per-user blocks are a subset of all unordered pairs among
$\sum_d \nleaf^{(d)}$ items, and the second by
$\sum_d \nleaf^{(d)} = \nleaf \leq \Nleaf$ and monotonicity of
$\binom{\cdot}{2}$. The verification-tag contribution satisfies
\[
  \sum_d \qv^{(d)} / 2^{\tauroot}
  \leq \Qv / 2^{\tauroot}.
\]
Combining with the $\mathsf{coll}$ penalty $D(D-1)/2^{k+1}$ and
$\Lift_c(\Ntot)$ yields the stated bound.
\end{proof}

\begin{theorem}[CMTDs-4 Committing Security]
\label{thm:cmtds4}
For every $s \geq 2$ and every $s$-way CMTDs-4
adversary $\mathcal{A}$ against \TW{} with the resources of
\Cref{sec:resources},
\[
  \Advcmtds{4}_\TW(\mathcal{A})
  \;\leq\;
  \Lift_c(\Ntot) + \RootColl_{\tauroot}(\Nprim, s)
  \;=\;
  \Lift_c(\Ntot) + \frac{\binom{\Nprim + s}{s}}{2^{\tauroot (s - 1)}}.
\]
\end{theorem}

\begin{proof}
For deterministic $\Dec$, every CMTDs-4 winner is
also a CMTDs-3 winner, yielding
\[
  \Advcmtds{4}_\TW(\mathcal{A}) \leq \Advcmtds{3}_\TW(\mathcal{A})
\]
by \Cref{app:committing-transfer}. The public-permutation
CMTDs-3 bound from \Cref{cor:cmtds3-pp},
\[
  \Advcmtds{3}_\TW(\mathcal{A})
  \leq \Lift_c(\Ntot) + \RootColl_{\tauroot}(\Nprim, s),
\]
then gives the claim.
\end{proof}

\begin{remark}[Looseness of the Root-Collision Cap]
\label{rem:rtag-cap}
The $\Nprim$ in \Cref{thm:cmtds4}'s root-collision term counts every
primitive query, not only those whose induced simulator $\ROflat$ call
lands at the root tag output point---$\Rtag$, or the elided-aggregation
$\Rmsg^+$ for single-chunk ciphertexts (\Cref{tab:keyed-twout}). By
\Cref{cor:output-sep}, calls
landing at any other output-point class cannot contribute to the
candidate sequence, and hence cannot contribute to the
candidate-collision event of \Cref{thm:cmtds3-ro}'s proof. A sharper
candidate-sequence analysis could replace $\Nprim$ in the root-collision
term by a class-filtered effective query count, no larger than $\Nprim$
and potentially smaller. The direct
CMTs-$\ell$ bounds of \Cref{app:committing-transfer} use the
same accounting boundary: their challenger encryption checks are
construction-internal work counted in $N$, so they do not enlarge the
direct-query cap $\Nprim$. The headline
CMTDs-4 statement nevertheless keeps the $\Nprim$ form for simplicity and
uniformity with \Cref{thm:ae,thm:mu-ind}; the birthday threshold at
$\tauroot = 256$ is unaffected.
\end{remark}

\begin{theorem}[Tag-Projected CMTDs-4 Committing Security]
\label{thm:tag-cmtds4}
For every $s \geq 2$ and every $s$-way tag-CMTDs-4
adversary $\mathcal{A}$ against \TW{} with the resources of
\Cref{sec:resources},
\[
  \Advtagcmtds{4}_\TW(\mathcal{A})
  \;\leq\;
  \Lift_c(\Ntot)
  + \LeafColl_{\tauleaf}(\Nprim+\Nleaf)
  + \RootColl_{\tauroot}(\Nprim, s).
\]
\end{theorem}

\begin{proof}
Immediate from the public-permutation corollary
\Cref{cor:tag-cmtds4-pp}, which lifts the random-oracle bound of
\Cref{thm:tag-cmtds4-ro}. The ordinary CMTDs-4 game is the special case
in which the tag-CMTDs-4 adversary uses the same ciphertext body in all
openings, but the theorem allows distinct bodies sharing only the final
root tag. The additional $\LeafColl$ term accounts for the only extra
way distinct bodies can merge before the root tag is compared: a
collision among the leaf tags absorbed by the root transcript.
\end{proof}

\begin{theorem}[Context Discoverability]
\label{thm:cdy}
For every context-discoverability adversary $\mathcal{A}$ against \TW{}
with the resources of \Cref{sec:resources},
\[
  \Advcdy_\TW(\mathcal{A})
  \;\leq\;
  \Lift_c(\Ntot) + \RootPre_{\tauroot}(\Nprim)
  \;=\;
  \Lift_c(\Ntot) + \frac{\Nprim + 1}{2^{\tauroot}}.
\]
\end{theorem}

\begin{proof}
Immediate from the public-permutation corollary \Cref{cor:cdy-pp},
which lifts the random-oracle bound of \Cref{thm:cdy-ro} through the
flat-sponge transform. Where \Cref{thm:cmtds4} bounds the probability
that $s$ distinct contexts collide on one root tag, this bounds the
probability that a single context reaches a root tag fixed in advance;
the candidate-collision argument of \Cref{thm:cmtds3-ro} is replaced by
the candidate-preimage argument of \Cref{thm:cdy-ro}, dropping the
exponent from $\tauroot(s-1)$ to $\tauroot$.
\end{proof}

\begin{theorem}[Tag-Projected Context Discoverability]
\label{thm:tag-cdy}
For every tag-projected context-discoverability adversary
$\mathcal{A}$ against \TW{} with the resources of
\Cref{sec:resources},
\[
  \Advtagcdy_\TW(\mathcal{A})
  =
  \Advcdy_\TW(\mathcal{A})
  \;\leq\;
  \Lift_c(\Ntot) + \RootPre_{\tauroot}(\Nprim).
\]
\end{theorem}

\begin{proof}
The CDY game of \Cref{sec:games-pp} fixes only the root tag $T^*$ and
lets the adversary choose the ciphertext body and context. This is
exactly the $\TagProj(C\concat T)=T$ projection version of
discoverability, so the equality of advantages is definitional and the
bound is \Cref{thm:cdy}.
\end{proof}

\subsection{Loss Terms}
\label{sec:loss-terms}

Five closed-form expressions recur across the headline bounds of
\Cref{sec:headline}---$\Lift_c$ throughout, with $\KeyGuess_k$,
$\RootColl_\tau$, $\RootPre_\tau$, and $\LeafColl_\tau$ in some. We
collect them here for reference.

\paragraph{Flat-sponge loss.}
The flat-sponge differentiating loss is
\[
  \Lift_c(\Ntot) = \Ntot(\Ntot + 1) / 2^{c + 1}.
\]
\cite[Lemmas~4 and~5]{BDPV08} give the differentiating probabilities
in exact product form: the random-transformation probability is
$f_T(\Ntot) = 1 - \prod_{i=1}^{\Ntot}(1 - i/2^c)$, and the factor of
the random-permutation probability $f_P(\Ntot)$ corresponding to the
$f_T$ factor at index $i$ is that factor divided by a quantity
$1 - (i-1)/2^{r+c} \in (0, 1]$, hence at least as large, so
$f_P(\Ntot) \leq f_T(\Ntot)$. For
$\Ntot < 2^c$ every factor lies in $[0, 1]$ and
\[
  f_P(\Ntot) \leq f_T(\Ntot)
  = 1 - \prod_{i=1}^{\Ntot}\Bigl(1 - \frac{i}{2^c}\Bigr)
  \leq \sum_{i=1}^{\Ntot} \frac{i}{2^c}
  = \frac{\Ntot(\Ntot + 1)}{2^{c+1}}
  = \Lift_c(\Ntot),
\]
the second inequality by the Weierstrass product inequality
$\prod_i (1 - x_i) \geq 1 - \sum_i x_i$ for $x_i \in [0, 1]$.
\cite[Theorem~2]{BDPV08} carries the side condition
$\Ntot < 2^c$; for
$\Ntot \geq 2^c$, $\Lift_c(\Ntot) \geq 1$ and the advantage bound is
trivial without any appeal to indifferentiability, and already at
$\Ntot \geq 2^{c/2}$ the bound exceeds $1/2$ and provides no
meaningful security.

\paragraph{Key-guessing term.}
The ideal-system key-guessing term of \cite{BDPV11} is
\[
  \KeyGuess_k(Q) = Q / 2^k.
\]
It accounts for direct $\ROflat$ queries whose inputs satisfy the
decoder $\KeyedTWOut$ of \Cref{tab:keyed-twout}. Each such query is an
adaptive test of the decoded candidate key against the hidden key $K$,
independent of the query's requested output length. \Cref{lem:key-tests}
bounds the probability that any of $Q$ direct queries tests the right
key by $Q / 2^k$.

\paragraph{Root-collision term.}
The $s$-way root tag multi-collision term is
\[
  \RootColl_\tau(Q, s) = \binom{Q + s}{s} \big/ 2^{\tau (s - 1)},
\]
with the abbreviation $\RootColl_{\tauroot}(Q, s) =
\RootColl_\tau(Q, s)$ at $\tau = \tauroot$. The $+s$ is candidate-sequence
slack for the $s$ final root transcript inputs generated by the
committing challenger after the adversary outputs its tuples;
these lookups are construction-internal and are not added to
$\qRO(\mathcal{B})$. In the public-permutation CMTDs-3 and
direct CMTs-$\ell$ bounds, the committing proof builds a
chronological candidate sequence and bounds only the event that some
$s$ distinct candidates have equal root tag projections. The flat-sponge
lift of
\Cref{app:flat-sponge-lift} gives $Q = \Nprim$: only the adversary's
$\Nprim$ direct primitive queries can induce simulator-side
$\ROflat$ queries that contribute to that sequence.

\paragraph{Root-preimage term.}
The root tag preimage term is
\[
  \RootPre_\tau(Q) = (Q + 1) \big/ 2^{\tau},
\]
with the abbreviation $\RootPre_{\tauroot}(Q) = \RootPre_\tau(Q)$ at
$\tau = \tauroot$. It is the preimage analogue of $\RootColl_\tau$: a
single query-independent target tag replaces the $s$-way internal
collision, the $+s$ candidate-sequence slack collapses to the $+1$ of
the single discoverability challenger decryption check, and the
exponent drops from $\tau(s-1)$ to $\tau$. As in the root-collision
term, the flat-sponge lift gives $Q = \Nprim$. The term bounds the
query-independent target of context discoverability (\Cref{thm:cdy}).

\paragraph{Leaf-collision term.}
The known-key leaf tag collision term is
\[
  \LeafColl_\tau(Q) = \binom{Q}{2} \big/ 2^\tau.
\]
It is the pairwise birthday bound on $\tau$-bit leaf tag projections
used by the tag-projected CMTDs-4 theorem. Unlike the hidden-key
leaf collision term in the INT-CTXT bound, the compact committing game
has adversary-supplied keys, so direct primitive queries can pre-read
candidate leaf tag points. The random-oracle proof therefore counts
both direct leaf tag candidates and challenger-generated final leaf
transcripts; after the flat-sponge lift this gives
$Q=\Nprim+\Nleaf$.

\subsection{Supporting Bounds}
\label{sec:thms-pp}

The headline theorems of \Cref{sec:headline} feed on five
random-oracle-model bounds---the privacy and INT-CTXT envelopes, the
CMTDs-3 bound, the tag-CMTDs-4 bound, and the CDY bound---whose proofs occupy
\Cref{sec:privacy,sec:authenticity,sec:cmtds3}; the remaining
committing-notion bounds live in \Cref{app:committing-transfer}.
Resource-indexed
suprema range over flat-RO adversaries (\Cref{sec:ro-model}) with the
displayed resources, and the flat-sponge lift of
\Cref{app:flat-sponge-lift} instantiates each envelope at $\qRO =
\Nprim$ when transferring to the public-permutation setting.

\paragraph{Privacy envelope.}
Define
\[
  \Epspriv(\QRO) = \sup\bigl\{
    \Adv^{\mathsf{priv}}_{\RO}(\mathcal{B}) :
    \qRO(\mathcal{B}) \leq \QRO
  \bigr\},
\]
the supremum taken over nonce-respecting privacy adversaries against
$\Enc_K^\RO$. The flat-RO privacy theorem of \Cref{sec:privacy} gives
\[
  \Epspriv(\QRO) \leq \KeyGuess_k(\QRO).
\]

\paragraph{INT-CTXT envelope.}
Define
\[
\begin{aligned}
  \Epsintctxt(\QRO, \Nleaf, \Qv)
  = \sup\bigl\{&
    \Adv^{\mathsf{int\text{-}ctxt}}_{\RO}(\mathcal{B}) : \\
  & \qRO(\mathcal{B}) \leq \QRO,\,
    \nleaf(\mathcal{B}) \leq \Nleaf,\,
    \qv(\mathcal{B}) \leq \Qv
  \bigr\}.
\end{aligned}
\]
The flat-RO INT-CTXT theorem of \Cref{sec:authenticity} gives
\[
  \Epsintctxt(\QRO, \Nleaf, \Qv)
  \leq \KeyGuess_k(\QRO)
   + \frac{\Nleaf(\Nleaf - 1)}{2^{\tauleaf + 1}}
   + \frac{\Qv}{2^{\tauroot}}.
\]

\paragraph{CMTDs-3 bound.}
The flat-RO CMTDs-3 theorem of \Cref{sec:cmtds3} is the
direct root-collision bound
\[
  \Advcmtds{3}_{\RO}(\mathcal{B})
  \leq \RootColl_{\tauroot}(\qRO(\mathcal{B}), s)
  = \binom{\qRO(\mathcal{B}) + s}{s} \big/ 2^{\tauroot (s - 1)},
\]
with no key-guessing or leaf-collision term. The keys are
adversary-supplied, so a direct query on an input accepted by the
keyed transcript decoder $\KeyedTWOut$ is one of the adversary's
visible $\ROflat$ queries and is counted in $\qRO(\mathcal{B})$;
leaf tag collisions cannot merge
the final root transcripts whose tags are compared, since by
\Cref{cor:triple-sep} distinct $(K, N, A)$ triples force distinct
final root transcripts before aggregation values matter.

\paragraph{Tag-CMTDs-4 bound.}
The flat-RO tag-CMTDs-4 theorem of \Cref{sec:cmtds3} is
\[
  \Advtagcmtds{4}_{\RO}(\mathcal{B})
  \leq
  \LeafColl_{\tauleaf}(\qRO(\mathcal{B})+\nleaf(\mathcal{B}))
  + \RootColl_{\tauroot}(\qRO(\mathcal{B}), s).
\]
The root-collision part is the same final-root candidate argument as
CMTDs-3. The additional known-key leaf-collision term is needed because
tag-projected openings may use different ciphertext bodies; absent such
a leaf tag collision, pairwise distinct full inputs force pairwise
distinct final root transcripts before the common tag is compared.

\paragraph{CDY bound.}
The flat-RO CDY theorem of \Cref{sec:cmtds3} is the direct
root-preimage bound
\[
  \Advcdy_{\RO}(\mathcal{B})
  \leq \RootPre_{\tauroot}(\qRO(\mathcal{B}))
  = (\qRO(\mathcal{B}) + 1) \big/ 2^{\tauroot},
\]
again with no key-guessing or leaf-collision term, by the same
adversary-supplied-key accounting. A single context replaces the $s$
tuples, so the candidate-collision argument becomes a
candidate-preimage argument against the query-independent target tag
and no triple-separation is invoked.
\label{sec:lift-sketch}

\Cref{app:flat-sponge-lift} carries out the flat-sponge lift in two
hops. Hop~1 applies the \cite{BDPV08} indifferentiability theorem
through \Cref{lem:lift-bound}, replacing the real permutation by the
corresponding flat-RO \TW{} construction interface and the \cite{BDPV08}
public-primitive simulator at cost at most $\Lift_c(\Ntot)$, where
$\Ntot = N + \Nprim$ counts construction-internal sponge calls and
direct primitive queries. Hop~2 (\Cref{lem:derived-adversary}) defines a
derived flat-RO adversary
$\mathcal{B} = \mathcal{B}_{\text{derived}}(\mathcal{A})$ that answers
each primitive query by running the simulator against direct $\ROflat$
queries; this is a lossless equality of joint distributions, and the
simulator-call filter of \Cref{lem:sim-ro-call-filter} makes each
primitive query induce at most one $\ROflat$ call, so
\[
  \qRO(\mathcal{B}_{\text{derived}}) \leq \Nprim.
\]
The random-oracle envelopes of \Cref{sec:thms-pp} are therefore
instantiated at $\Nprim$, not $\Ntot$, in the headline theorems---as is
the CMTDs-3 candidate-sequence length bound $m \leq \Nprim + s$ of
\Cref{sec:cmtds3}, the tag-CMTDs-4 leaf/root candidate bounds, and the
direct CMTs-$\ell$ bounds of
\Cref{app:committing-transfer}, whose challenger encryption checks are
construction-internal and counted in $N$, not in $\Nprim$.

\subsection{Concrete TW128 Bounds}
\label{sec:concrete}

The implemented parameters of \Cref{sec:params} fix
\[
  k = c = \tauroot = \tauleaf = 256.
\]
Substituting into
\Cref{thm:ae,thm:mu-ind,thm:cmtds4,thm:tag-cmtds4,thm:cdy,thm:tag-cdy}
and using the random-oracle envelopes of \Cref{sec:thms-pp} gives the
following closed-form bounds.

\paragraph{All-in-one AE.}
\[
  \Advae_\TW(\mathcal{A})
  \;\leq\;
  \frac{\Ntot(\Ntot + 1)}{2^{257}}
  + \frac{2 \Nprim + \Qv}{2^{256}}
  + \frac{\Nleaf(\Nleaf - 1)}{2^{257}}.
\]

\paragraph{Multi-user indistinguishability.}
\[
  \Advmuind_\TW(\mathcal{A})
  \;\leq\;
  \frac{\Ntot(\Ntot + 1)}{2^{257}}
  + \frac{2 D \Nprim + \Qv}{2^{256}}
  + \frac{D(D - 1) + \Nleaf(\Nleaf - 1)}{2^{257}}.
\]

\paragraph{$s$-way CMTDs-4 committing security.}
For every $s \geq 2$,
\[
  \Advcmtds{4}_\TW(\mathcal{A})
  \;\leq\;
  \frac{\Ntot(\Ntot + 1)}{2^{257}}
  + \frac{\binom{\Nprim + s}{s}}{2^{256 (s - 1)}}.
\]

\begin{remark}[Two-Key CMTD-4 Specialization]
\label{rem:cmtd4-s2}
Setting $s = 2$ recovers the ordinary two-key
CMTD-4 statement of \cite{BH22}:
\[
  \Advcmtds{4}_\TW(\mathcal{A})
  \;\leq\;
  \frac{\Ntot(\Ntot + 1)}{2^{257}}
  + \frac{\binom{\Nprim + 2}{2}}{2^{256}}
  \;=\;
  \frac{\Ntot(\Ntot + 1) + (\Nprim + 1)(\Nprim + 2)}{2^{257}}.
\]
\end{remark}

\paragraph{$s$-way tag-projected CMTDs-4 committing security.}
For every $s \geq 2$,
\[
  \Advtagcmtds{4}_\TW(\mathcal{A})
  \;\leq\;
  \frac{\Ntot(\Ntot + 1)}{2^{257}}
  + \frac{\binom{\Nprim+\Nleaf}{2}}{2^{256}}
  + \frac{\binom{\Nprim + s}{s}}{2^{256 (s - 1)}}.
\]
For $s=2$, this is the compact two-opening bound in which the common
commitment is only the tag:
\[
\begin{aligned}
  \Advtagcmtds{4}_\TW(\mathcal{A})
\leq\;&
  \frac{\Ntot(\Ntot + 1)}{2^{257}}
  + \frac{(\Nprim+\Nleaf)(\Nprim+\Nleaf-1)}{2^{257}} \\
& + \frac{(\Nprim + 1)(\Nprim + 2)}{2^{257}} .
\end{aligned}
\]

\paragraph{Context discoverability.}
\[
  \Advtagcdy_\TW(\mathcal{A})
  =
  \Advcdy_\TW(\mathcal{A})
  \;\leq\;
  \frac{\Ntot(\Ntot + 1)}{2^{257}}
  + \frac{\Nprim + 1}{2^{256}}.
\]

\subsection{Concrete Security Claim}
\label{sec:concrete-claim}

\TW{} attains a $128$-bit security target across the standard notions of
\Cref{sec:concrete}. The birthday term $\Lift_c(\Ntot) = \Ntot(\Ntot +
1) / 2^{257}$ dominates the AE bound and is a leading term in the
committing bounds, where the two-key specialization also has a
co-leading root-collision term; it contributes alongside a multi-user
hybrid term in the mu-ind bound; and it is the sole leading term in the
discoverability bound, whose root-preimage term is linear in $\Nprim$.
The tag-projected CMTDs-4 strengthening adds one leaf tag birthday term on
$\Nprim+\Nleaf$. The headline work cap
$\Ntot \leq 2^{63}$ keeps the resulting total below $2^{-128}$ after
the resource consequences of
\Cref{sec:resources} are applied. Namely, $\Nprim \leq \Ntot$ by
definition, and the pathwise bounds $\qv \leq N$ and $\nleaf \leq N$
allow the envelope parameters $\Qv$ and $\Nleaf$ to be instantiated at
values at most $\Ntot$. The mu-ind user cap $D$ remains an independent hypothesis,
since created-but-idle users need not contribute to $N$.
\begin{itemize}
\item \emph{AE.} For every nonce-respecting AE adversary against \TW{}
  with $\Ntot \leq 2^{63}$ \Keccakp{} calls,
  $\Advae_\TW(\mathcal{A}) \leq 2^{-128}$. The flat-sponge and
  leaf-collision terms are each bounded by one $2^{-131}$-sized
  birthday term, up to lower-order summands, while
  $(2\Nprim+\Qv)/2^{256} \leq 3 \cdot 2^{-193}$.
\item \emph{Multi-user indistinguishability.} For every $D$-user
  mu-ind adversary against \TW{} with $\Ntot \leq 2^{63}$ \Keccakp{}
  calls and $D \leq 2^{63}$ users, $\Advmuind_\TW(\mathcal{A}) \leq
  2^{-128}$. At the threshold, the hybrid term satisfies
  $2D\Nprim/2^{256} \leq 2^{-129} = 4 \cdot 2^{-131}$; the
  flat-sponge term is at most $2^{-131}+2^{-194}$; the user-key
  birthday and leaf-collision terms are each below $2^{-131}$; and
  $\Qv/2^{256} \leq 2^{-193}$. Their sum is at most
  $7 \cdot 2^{-131} + 3 \cdot 2^{-194} < 8 \cdot 2^{-131} = 2^{-128}$.
\item \emph{Committing security.} For every $s \geq 2$ and every
  $s$-way CMTDs-4 adversary against \TW{} with $\Ntot \leq
  2^{63}$ \Keccakp{} calls (so in particular $\Nprim \leq 2^{63}$),
  $\Advcmtds{4}_\TW(\mathcal{A}) \leq 2^{-128}$. In the two-key case
  $s=2$, the flat-sponge and root-collision terms are each bounded by
  one $2^{-131}$-sized birthday term, up to lower-order summands. For
  $s \geq 3$, the root-collision term decays as $2^{-256(s - 1)}$ in
  the denominator, so it becomes negligible relative to the birthday
  term well before $\Ntot$ saturates.
\item \emph{Tag-projected committing security.} For every $s \geq 2$
  and every $s$-way tag-CMTDs-4 adversary against \TW{} with $\Ntot
  \leq 2^{63}$ \Keccakp{} calls, $\Advtagcmtds{4}_\TW(\mathcal{A})
  \leq 2^{-128}$. At $s=2$, the flat-sponge and root-collision terms
  are each at most one $2^{-131}$-sized birthday term, while
  $\LeafColl_{256}(\Nprim+\Nleaf)$ is at most
  $\binom{2^{64}}{2}/2^{256} < 2^{-129}$ using
  $\Nprim,\Nleaf \leq \Ntot$. Larger $s$ only decreases the
  root-collision contribution.
\item \emph{Context discoverability.} For every CDY adversary against
  \TW{} with $\Ntot \leq 2^{63}$ \Keccakp{} calls (so in particular
  $\Nprim \leq 2^{63}$),
  $\Advtagcdy_\TW(\mathcal{A})=\Advcdy_\TW(\mathcal{A}) \leq
  2^{-128}$. The
  flat-sponge term is one $2^{-131}$-sized birthday term, up to
  lower-order summands, and the root-preimage term satisfies
  $(\Nprim + 1)/2^{256} \leq 2 \cdot 2^{-193}$, far below the
  $128$-bit target; as in the committing bound, the flat-sponge
  differentiating loss dominates.
\end{itemize}

Expressed in bytes of construction work, $\Ntot \leq 2^{63}$
corresponds to processing at most $\rho \cdot 2^{63} = 167 \cdot
2^{63} < 2^{71}$ bytes of plaintext, associated data, ciphertext, and
aggregation through the duplex (\Cref{sec:params}), well beyond any
realistic deployment lifetime. The estimates above are direct
instantiations of the closed-form expressions of \Cref{sec:concrete}
using the derived resource bounds from \Cref{sec:resources}; no separate
top-level cap on $\Qv$ or $\Nleaf$ is required.

\subsection{Discussion of the Bounds}
\label{sec:bound-discussion}

The closed-form bounds of \Cref{sec:concrete} aggregate four loss
terms: the flat-sponge differentiating loss $\Lift_c$, the
key-guessing contribution $\KeyGuess_k$ (entering privacy and
INT-CTXT independently, hence the factor of two in the AE and
mu-ind bounds), the leaf tag pairwise-collision term, and the
verification-tag term. The CMTDs-4 bound substitutes the
$s$-way root-collision term $\RootColl_{\tauroot}$ for the latter
three; the tag-projected CMTDs-4 bound adds the known-key
$\LeafColl_{\tauleaf}$ term; and the CDY bound uses the root-preimage
term $\RootPre_{\tauroot}$.
We comment on dominance and tightness.

\paragraph{Dominance.}
In the single-user AE bound, the birthday-on-capacity loss
$\Lift_c(\Ntot) = \Ntot(\Ntot + 1) / 2^{257}$ dominates for
$6 \leq \Ntot < 2^{c}$. The key-guessing and
verification-tag contributions $(2 \Nprim + \Qv) / 2^{256}$ are
linear in their resource caps and bounded by $3 \Ntot / 2^{256}$,
hence below $\Lift_c$ for $\Ntot \geq 6$. The leaf tag collision
term $\Nleaf(\Nleaf - 1) / 2^{257}$ is at most $\Lift_c$ via
$\Nleaf \leq \Ntot$. The CMTDs-4 bound is similar: the
$s = 2$ root-collision term equals $(\Nprim + 1)(\Nprim + 2) /
2^{257}$, which is within a small additive $\Ntot$-term of
$\Lift_c$ under $\Nprim \leq \Ntot$, and it decays as $2^{-256(s -
1)}$ in the denominator for $s \geq 3$. The tag-projected CMTDs-4
bound has the same root-collision term plus
$\LeafColl_{256}(\Nprim+\Nleaf)$, which is at most
$\binom{2\Ntot}{2}/2^{256}$ under the resource consequences of
\Cref{sec:resources}; at the concrete cap this is below $2^{-129}$.
The CDY bound is dominated by
$\Lift_c$ outright: its root-preimage term $(\Nprim + 1) / 2^{256}$ is
linear in $\Nprim \leq \Ntot$ and hence below $\Lift_c$ for
$\Ntot \geq 3$.

The multi-user mu-ind bound is the exception. For
$D \leq \Ntot(\Ntot+1) / (4 \Nprim)$ users (typical for any
practical deployment), the hybrid-induced term remains at most
$\Lift_c$. At the upper end of \Cref{sec:concrete-claim}'s parameters
($\Ntot, D, \Nprim \leq 2^{63}$), the hybrid-induced term
$2 D \cdot \KeyGuess_k(\Nprim) \leq 2 D \Ntot / 2^{256}$ reaches
$\approx 4 \Lift_c$ and is the dominant term; the user-key
birthday $D(D - 1) / 2^{257}$ and the leaf tag collision each
contribute roughly $\Lift_c$. The $2^{-128}$ figure of
\Cref{sec:concrete-claim} is an upper bound on the sum of these terms
at the threshold.

\paragraph{Tightness.}
\begin{itemize}
\item \emph{$\KeyGuess_k(Q) = Q / 2^k$.} Tight at constants: a direct
  random-guess attack on a $k$-bit key matches the per-query success
  probability $1 / 2^k$.
\item \emph{$\RootColl_\tau(Q, s) = \binom{Q + s}{s} / 2^{\tau (s -
  1)}$.} Tight up to lower-order terms: it is the standard $s$-way
  multi-collision bound on a $\tau$-bit projection of $\ROflat$, and
  a generic birthday-style adversary against the random oracle
  matches the leading-order term.
\item \emph{Leaf tag collision $\Nleaf (\Nleaf - 1) / 2^{\tauleaf +
  1}$.} The pairwise union bound is tight up to constants: a generic
  birthday attack on the $\tauleaf$-bit leaf tag projection matches
  the leading $\Nleaf^2 / 2^{\tauleaf + 1}$ rate.
\item \emph{Known-key leaf collision $\LeafColl_\tau(Q)$.} The same
  birthday attack applies when keys are adversary-supplied and direct
  queries can target leaf tag points; the compact CMTDs-4 proof uses
  $Q=\Nprim+\Nleaf$.
\item \emph{Flat-sponge lift.} The term
  $\Lift_c(\Ntot) = \Ntot(\Ntot + 1) / 2^{c + 1}$ bounds the exact
  random-permutation differentiating probability of
  \cite[Lemma~5]{BDPV08} via \Cref{sec:loss-terms}. The
  remaining slack is the unimproved
  tightness of the \cite{BDPV08} simulator itself, which is an open
  problem on the sponge side and not specific to \TW{}.
\end{itemize}

Cross-scheme bound-level comparisons against prior sponge AEADs and
the committing-AEAD line of \cite{BH22} are deferred to
\Cref{sec:compare-table,sec:compare-sponge,sec:compare-commit}.
